\documentclass[aspectratio=169, table, xcdraw]{beamer}

\usetheme{Madrid}
\usecolortheme{default}
\definecolor{ucbblue}{RGB}{0, 51, 102}

\setbeamercolor{structure}{fg=ucbblue}
\setbeamercolor*{palette primary}{fg=white,bg=ucbblue}
\setbeamercolor*{palette secondary}{fg=white,bg=ucbblue!80}
\setbeamercolor*{palette tertiary}{fg=white,bg=ucbblue!90}
\setbeamercolor{title}{fg=white, bg=ucbblue}
\setbeamercolor{frametitle}{fg=white, bg=ucbblue}
\setbeamercolor{item}{fg=ucbblue}

\usepackage[T1]{fontenc}
\usepackage[utf8]{inputenc}
\usepackage{tgpagella}
\usefonttheme{serif}
\usepackage[spanish, es-noshorthands]{babel}
\usepackage{amsmath, amssymb, bm}
\usepackage{graphicx}
\usepackage[most]{tcolorbox}
\usepackage{booktabs}
\usepackage{tabularx}
\usepackage{tikz}
\usepackage{pgfplots}
\pgfplotsset{compat=1.18}
\usepackage[american]{circuitikz}
\usetikzlibrary{shapes.geometric, arrows.meta, positioning, calc}

\title[Interacción y Multietapa]{Consolidación de emisor común, interacción fuente-amplificador-carga y puente hacia multietapa}
\subtitle{IMT-221 Circuitos Electrónicos III — Semana 6, Sesión 3}
\author[B. Quiroga Turdera]{M.Sc. Ing. Bernardo Quiroga Turdera}
\institute[UCB Tarija]{Universidad Católica Boliviana ``San Pablo'' — Sede Tarija}
\date{Semana 6}

\begin{document}

\begin{frame}
    \titlepage
\end{frame}

\begin{frame}{El Amplificador Completo}
    \begin{tcolorbox}[colback=ucbblue!5, colframe=ucbblue, title=Idea Central de la Sesión]
        El amplificador útil es el sistema completo \textbf{fuente–transistor–carga}, no la ecuación del transistor en aislamiento.
    \end{tcolorbox}
    
    \vspace{0.4cm}
    \textbf{Cadena de Análisis:}
    \begin{center}
        $Q \rightarrow g_m, r_\pi \rightarrow \text{Híbrido-}\pi \rightarrow A_v \rightarrow R_{in} \rightarrow R_S \rightarrow R_L \rightarrow R_{out} \rightarrow G_v$
    \end{center}
\end{frame}

\begin{frame}{Diagnóstico Breve: Emisor Común Aislado}
    \begin{columns}
        \begin{column}{0.4\textwidth}
            \centering
            \begin{circuitikz}[scale=0.9, transform shape, thick]
                \draw (0,0) node[npn](Q){};
                \draw (Q.E) node[ground]{};
                \draw (Q.C) to[R, l_={$R_C$}] (0,2) -- (0,2.5) node[above]{$+V_{CC}$};
                \draw (Q.B) -- (-2,0) node[left]{$v_i$};
                \draw (Q.C) to[short, -o] (1,0.77) node[right]{$v_o$};
            \end{circuitikz}
        \end{column}
        \begin{column}{0.55\textwidth}
            \textbf{Condiciones de Operación (Asignadas):}
            \begin{itemize}
                \item $I_{CQ} = 1.0\,\text{mA}$
                \item $\beta = 100$
                \item $R_C = 4.7\,\text{k}\Omega$
                \item $T \approx 300\,\text{K} \implies V_T \approx 25.85\,\text{mV}$
                \item Supuestos: Activa directa, $r_o \to \infty$, $v_i = v_\pi$
            \end{itemize}
            
            \vspace{0.2cm}
            \textbf{Parámetros Mínimos:}
            \begin{align*}
                g_m &= \frac{1\,\text{mA}}{25.85\,\text{mV}} \approx 38.68\,\text{mS} \\
                r_\pi &= \frac{100}{38.68\,\text{mS}} \approx 2.585\,\text{k}\Omega
            \end{align*}
        \end{column}
    \end{columns}
\end{frame}

\begin{frame}{Atenuación de Entrada (Efecto de $R_S$)}
    Una fuente real tiene resistencia interna ($R_S$). La red de polarización ($R_B$) y el transistor ($r_\pi$) forman una resistencia de entrada que "carga" a la fuente.
    
    \vspace{0.2cm}
    \begin{columns}
        \begin{column}{0.55\textwidth}
            \centering
            \begin{circuitikz}[scale=0.85, transform shape, thick]
                \draw (0,0) to[sV, l={$v_s$}] (0,2) to[R, l={$R_S$}] (2,2);
                \draw (2,0) to[R, l={$R_{in}$}] (2,2);
                \draw (2,2) to[short, -o] (3,2) node[right]{$v_i$};
                \draw (0,0) -- (2,0) node[ground]{};
            \end{circuitikz}
        \end{column}
        \begin{column}{0.4\textwidth}
            \textbf{Resistencia de entrada:}
            \begin{equation*}
                R_{in} = R_B \parallel r_\pi
            \end{equation*}
            
            \vspace{0.2cm}
            \textbf{Divisor de tensión:}
            \begin{equation*}
                \frac{v_i}{v_s} = \frac{R_{in}}{R_S + R_{in}} < 1
            \end{equation*}
        \end{column}
    \end{columns}
    \vspace{0.3cm}
    \textit{Si $R_S$ es comparable o mayor que $R_{in}$, el amplificador recibe sólo una fracción de la señal de la fuente.}
\end{frame}

\begin{frame}{Efecto de la Carga Externa ($R_L$)}
    La carga externa se acopla en AC al nodo del colector, quedando en paralelo con $R_C$.
    
    \vspace{0.2cm}
    \begin{columns}
        \begin{column}{0.45\textwidth}
            \centering
            \textbf{Colector Sin Carga} \\
            \vspace{0.2cm}
            \begin{circuitikz}[scale=0.8, transform shape, thick]
                \draw (0,2) node[ground, rotate=180]{} to[R, l_={$R_C$}] (0,0) node[below]{C};
                \draw (0,0) to[short, -o] (1.5,0) node[right]{$v_o$};
            \end{circuitikz}
            \vspace{0.2cm}
            \begin{equation*}
                A_{v0} = -g_m R_C
            \end{equation*}
        \end{column}
        \begin{column}{0.1\textwidth}
            $\implies$
        \end{column}
        \begin{column}{0.45\textwidth}
            \centering
            \textbf{Colector Con Carga} \\
            \vspace{0.2cm}
            \begin{circuitikz}[scale=0.8, transform shape, thick]
                \draw (0,2) node[ground, rotate=180]{} to[R, l_={$R_C$}] (0,0) node[below]{C};
                \draw (0,0) -- (1.5,0) to[R, l={$R_L$}] (1.5,2) node[ground, rotate=180]{};
                \draw (0,0) to[short, *-o] (2.5,0) node[right]{$v_o$};
            \end{circuitikz}
            \vspace{0.2cm}
            \begin{equation*}
                A_v = -g_m (R_C \parallel R_L)
            \end{equation*}
        \end{column}
    \end{columns}
    \vspace{0.3cm}
    \textit{Nota: $R_{out} \approx R_C$ (asumiendo $r_o \to \infty$). Agregar $R_L$ invariablemente disminuye la magnitud de la ganancia del nodo.}
\end{frame}

\begin{frame}{Jerarquía de Ganancias: $A_{v0}$, $A_v$, $G_v$}
    Definamos claramente las tres ganancias clave:
    
    \vspace{0.3cm}
    \begin{itemize}
        \item \textbf{$A_{v0} = -g_m R_C$:} Ganancia ideal intrínseca de la etapa sin carga externa.
        \item \textbf{$A_v = -g_m (R_C \parallel R_L)$:} Ganancia de la etapa desde su entrada $v_i$ hasta su salida $v_o$, afectada por $R_L$.
        \item \textbf{$G_v = \frac{v_o}{v_s} = \left(\frac{v_i}{v_s}\right) A_v$:} Ganancia \textbf{global} desde el generador de la fuente hasta la salida, afectada por $R_S$ y $R_L$.
    \end{itemize}
    
    \begin{alertblock}{Conclusión Crítica}
        Para predecir una medición real en laboratorio, \textbf{$G_v$} es el parámetro que describe el sistema completo.
    \end{alertblock}
\end{frame}

\begin{frame}{Ejemplo Integrado}
    \textbf{Datos:} $I_{CQ}=1\,\text{mA}$ ($g_m=38.68\,\text{mS}$, $r_\pi=2.585\,\text{k}\Omega$), $R_B=100\,\text{k}\Omega$, $R_S=10\,\text{k}\Omega$, $R_C=4.7\,\text{k}\Omega$, $R_L=10\,\text{k}\Omega$.
    
    \begin{columns}
        \begin{column}{0.4\textwidth}
            \resizebox{\textwidth}{!}{
            \begin{circuitikz}[thick]
                % Source part
                \draw (-2,0) to[sV, l={$v_s$}] (-2,2) to[R, l={$R_S$}] (0,2);
                \draw (-2,0) -- (0,0) node[ground]{};
                
                % Ampli Input
                \draw (0,2) to[short, *-] (2,2) node[right]{B};
                \draw (0,2) to[R, l_={$R_B$}] (0,0);
                \draw (2,2) to[R, l={$r_\pi$}] (2,0);
                \draw (0,0) -- (2,0) node[ground]{};
                
                % Ampli Output
                \draw (5,2) node[left]{C} to[cI, l_={$g_m v_\pi$}] (5,0) node[ground]{};
                \draw (5,2) -- (7,2);
                \draw (6,2) to[R, l={$R_C$}] (6,0) node[ground]{};
                \draw (7,2) to[R, l={$R_L$}] (7,0) node[ground]{};
                \draw (7,2) to[short, *-o] (8,2) node[right]{$v_o$};
            \end{circuitikz}
            }
        \end{column}
        \begin{column}{0.55\textwidth}
            \begin{enumerate}
                \item $R_{in} = 100\text{k} \parallel 2.585\text{k} \approx 2.520\,\text{k}\Omega$
                \item $\frac{v_i}{v_s} = \frac{2.520}{10 + 2.520} \approx 0.2013$
                \item $R_C' = 4.7\text{k} \parallel 10\text{k} \approx 3.197\,\text{k}\Omega$
                \item $A_v = -g_m R_C' \approx -123.7$
                \item $G_v = (0.2013)(-123.7) \approx -24.9$
            \end{enumerate}
        \end{column}
    \end{columns}
\end{frame}

\begin{frame}{Comparación Funcional: CE, CC, CB}
    \centering
    \renewcommand{\arraystretch}{1.4}
    \begin{tabularx}{0.95\textwidth}{|X|X|X|X|}
    \hline
    \rowcolor{ucbblue!20}
    \textbf{Propiedad} & \textbf{Emisor Común (CE)} & \textbf{Colector Común (CC)} & \textbf{Base Común (CB)} \\ \hline
    \textbf{Entrada $\to$ Salida} & Base $\to$ Colector & Base $\to$ Emisor & Emisor $\to$ Colector \\ \hline
    \textbf{Terminal Común} & Emisor & Colector & Base \\ \hline
    \textbf{Ganancia Voltaje} & Magnitud alta ($>1$) & Cercana a 1 ($\approx 1$) & Magnitud alta ($>1$) \\ \hline
    \textbf{Fase} & Invertida ($180^\circ$) & No invertida & No invertida \\ \hline
    \textbf{Z de Entrada ($R_{in}$)} & Moderada & \textbf{Alta} & \textbf{Baja} \\ \hline
    \textbf{Z de Salida ($R_{out}$)} & Moderada/Alta & \textbf{Baja} & Alta \\ \hline
    \textbf{Rol Típico} & Amplificador principal de tensión & \textbf{Buffer} (adaptador de impedancia) & Interfaz de corriente / Alta frecuencia \\ \hline
    \end{tabularx}
\end{frame}

\begin{frame}{Carga Interetapa y el Rol del Buffer (CC)}
    En una cascada (múltiples etapas), la entrada de la Etapa 2 actúa como la carga ($R_L$) de la Etapa 1.
    
    \vspace{0.3cm}
    \begin{columns}
        \begin{column}{0.4\textwidth}
            \centering
            \begin{circuitikz}[scale=0.8, transform shape, thick]
                \draw (0,0) to[sV, l={$v_{th1}$}] (0,2) to[R, l={$R_{out1}$}] (2,2);
                \draw (2,2) to[R, l={$R_{in2}$}] (2,0);
                \draw (0,0) -- (2,0) node[ground]{};
                \draw (2,2) to[short, *-o] (3,2) node[right]{$v_2$};
            \end{circuitikz}
        \end{column}
        \begin{column}{0.55\textwidth}
            La Etapa 1 entrega a la Etapa 2:
            \begin{equation*}
                v_2 = v_{th1} \left( \frac{R_{in2}}{R_{out1} + R_{in2}} \right)
            \end{equation*}
            Si $R_{in2} \ll R_{out1}$, la atenuación destruye la ganancia previa. \\
            \vspace{0.2cm}
            \textbf{Solución:} Insertar un \textbf{Colector Común (Buffer)} entre la ganancia y la carga pesada (su $R_{in}$ alto protege a la primera etapa, su $R_{out}$ bajo maneja bien la carga final).
        \end{column}
    \end{columns}
\end{frame}

\begin{frame}{Puente Cualitativo al Par Diferencial (Hito 2)}
    El Par Diferencial no son "dos Emisores Comunes independientes"; es una estructura acoplada que comparte polarización.
    
    \begin{columns}
        \begin{column}{0.5\textwidth}
            \centering
            \begin{circuitikz}[scale=0.7, transform shape, thick]
                % Q1 and Q2
                \draw (0,0) node[npn](Q1){}; \node[right] at (Q1) {Q1};
                \draw (4,0) node[npn, xscale=-1](Q2){}; \node[left] at (Q2) {Q2};
                % Bias
                \draw (Q1.C) to[R, l={$R_C$}] (0,3) -- (4,3) to[R, l_={$R_C$}] (Q2.C);
                \draw (2,3) -- (2,3.5) node[above]{$+V_{CC}$};
                \draw (Q1.E) -- (0,-1) -- (4,-1) -- (Q2.E);
                \draw (2,-1) to[I, l={$I_{TAIL}$}] (2,-3) node[below]{$-V_{EE}$};
                % Inputs
                \draw (Q1.B) to[short, -o] (-1.5,0) node[left]{$v_1$};
                \draw (Q2.B) to[short, -o] (5.5,0) node[right]{$v_2$};
                % Outputs
                \draw (Q1.C) to[short, *-o] (-1,2) node[left]{$v_{o1}$};
                \draw (Q2.C) to[short, *-o] (5,2) node[right]{$v_{o2}$};
            \end{circuitikz}
        \end{column}
        \begin{column}{0.48\textwidth}
            \textbf{Sensibilidad a la diferencia ($v_{id} = v_1 - v_2$):}
            \begin{itemize}
                \item Si $v_1 = v_2$, la corriente $I_{TAIL}$ se divide simétricamente.
                \item Si $v_1 > v_2$, Q1 conduce más corriente, robándosela a Q2. 
                \item El voltaje de colector de Q1 cae, el de Q2 sube.
            \end{itemize}
            \vspace{0.2cm}
            \textit{Comprender $g_m$ y la simetría es vital para medir pequeñas caídas en shunts de corriente.}
        \end{column}
    \end{columns}
\end{frame}

\end{document}
